\begin{problem}
  Translate the rectangular coordinates $(-5, 2)$ into polar coordinates $(r,
  \theta)$. (Approximations are acceptable here but they must be denoted
  correctly.) \textbf{[$5$ points]}
\end{problem}

\begin{solution}
  First, I need to get the radius ($r$):

  \begin{align*}
    r &= \sqrt{x^{2} + y^{2}} \\
      &= \sqrt{-5^{2} + 2^{2}} \\
      &= \sqrt{25 + 4} \\
    r &= \sqrt{29} \\
  .\end{align*}

  Now, I need to find $\theta$:

  \begin{align*}
    \theta = \tan^{-1} \left(-\frac{2}{5}\right) \\
    \theta = -0.38^{\circ} + 180^{\circ} = 179.62^{\circ}
  .\end{align*}

  So, my final answer is:
  \[ (\sqrt{29}, 179.62^{\circ}) \].
\end{solution}

\newpage

\begin{problem}
  Translate the polar coordinates $(5, \frac{7\pi}{6})$ into rectangular
  coordinates $(x,y )$. Provide \textbf{exact} values. \textbf{[$5$ points]}
\end{problem}

\begin{solution}
  \begin{align*}
    P &= (r\cos (\theta), r\sin (\theta)) \\
      &= (5\cos \left(\frac{7\pi}{6}\right), 5\sin \left(\frac{7\pi}{6}\right)) \\
      &= \left(-\frac{5\sqrt{3}}{2}, -\frac{5}{2}\right)
  .\end{align*}
\end{solution}

\newpage

\begin{problem}
  The polar ordered pair $(8, \frac{5\pi}{4})$ can be plotted as a "dot" on the
  polar coordinate plane. List two other (different) polar ordered pairs that
  are plotted on the same "dot". \textbf{[$5$ points]}
\end{problem}

\begin{solution}
  The first polar ordered pair that can be plotted as a "dot" on the polar
  coordinate plane is $(-8, \frac{5\pi}{4})$. Another coordinate is $(8,
  \frac{\pi}{4})$.
\end{solution}

\newpage

\begin{problem}
  Find a polar form, $z = re^{i\theta}$, of the complex number $z = -6 +
  2\sqrt{3}i$. Use exact values. \textbf{[$5$ points]}
\end{problem}

\begin{solution}
  First, let's find $r$:

  \begin{align*}
    r &= \sqrt{(-6)^{2} + (2\sqrt{3})^{2}} \\
      &= \sqrt{36 + 12} \\
      &= \sqrt{48} \\
      &= 4\sqrt{3} \\
  .\end{align*}

  Now, let's find $\theta$:

  \begin{align*}
    \qquad&\tan (\theta) = -\frac{2\sqrt{3}}{6} \\
    \implies\qquad&\theta = -\tan^{-1} (\frac{2\sqrt{3}}{6}) \\
    \implies\qquad&\theta = -\frac{\pi}{6} \\
  .\end{align*}

  The complex number $z = -6 + 2\sqrt{3}i$ can be represented in polar form:

  \[ z = 4\sqrt{3}e^{i\left(-\frac{\pi}{6}\right)} \].
\end{solution}
